The Length-Width Ratio (LW) provides a direct measure of district elongation:
a district with LW $= 4$ is four times as long as it is wide,
while a perfectly square district has LW $= 1$.
Unlike perimeter-based metrics (PP, S) that capture overall boundary irregularity
and area-based metrics (Reock, CH) that capture departure from convexity or
circularity, LW specifically measures the aspect ratio of the smallest rectangle
that contains the district.

\subsection{Intuitive Appeal}

LW's intuitive appeal in court proceedings is that a $3:1$ aspect ratio
(LW $= 3$) is easy to visualise: a district that is three times as long as
it is wide appears obviously elongated to any observer.
Courts do not need expert testimony to understand that a district that spans
the width of a state as a thin corridor is elongated; the LW $> 3$ threshold
formalises this visual intuition into a numerical criterion.

\subsection{Orientation Dependence and the AABB Problem}

The naive elongation measure---the axis-aligned bounding box (AABB)---reports
the width and height of the smallest north-south and east-west rectangle
containing the district.
AABB is computationally trivial (a single pass over vertex coordinates) but
is fundamentally orientation-dependent: a $1 \times 4$ rectangle oriented
at 45 degrees has AABB dimensions $\sqrt{2} \times \sqrt{2}(1 + 4/\sqrt{2}) \approx 3.54 \times 3.54$,
giving AABB LW $\approx 1.0$---dramatically underestimating the true elongation.

More critically, for a $1 \times 3$ rectangle oriented at 45 degrees:
width $= $ height $= (1 + 3)/\sqrt{2} \cdot \cos(45) = 4/\sqrt{2} \approx 2.83$.
AABB reports LW $= 2.83 / 2.83 = 1.0$, while the true LW $= 3$.
Opposing counsel could use AABB to argue that an elongated diagonal district
is ``equidimensional.''

The minimum bounding rectangle (MBR)---the smallest-area rectangle at any
orientation---correctly reports LW $= 3$ regardless of district orientation.
MBR is rotation-invariant; AABB is not.

\subsection{Scope}

This paper documents the rotating calipers algorithm for MBR computation,
the LW $> 3$ court threshold, and bisect's performance on the LW metric
across NC, WI, and TX.
